

\newcommand{\tnote}[1]{\TM{ [#1] }}

\section{Introduction} \label{sec:intro}
Training neural networks requires optimizing non-convex losses, which is often practically feasible but still not theoretically understood. The lack of understanding of non-convex optimization also limits the design of new principled optimizers for training neural networks that use theoretical insights. 

Early analysis on optimizing neural networks with linear or quadratic activations~\citep{du2018power, gunasekar2018implicit, li2018algorithmic,gunasekar2018characterizing,soltanolkotabi2018theoretical,hardt2016identity,kawaguchi2016deep,hardt2018gradient,oymak2020toward} relies on linear algebraic tools that do not extend to nonlinear and non-quadratic activations. The neural tangent kernel (NTK) approach analyzes nonconvex optimization under certain hyperparameter settings, e.g., when the initialization scale is large and the learning rate is small (see, e.g., ~\citet{du2018gradient,jacot2018neural, li2018learning, arora2019fine, daniely2016toward}). 
However, subsequent research shows that neural networks trained with practical hyperparameter settings typically outperform their corresponding NTK kernels~\citep{arora2019exact}. Furthermore, the initialization and learning rate under the NTK regime does not yield optimal generalization guarantees~\citep {wei2019regularization,chizat2018note,woodworth2020kernel, ghorbani2020neural}. 

Many recent works study modified versions of stochastic gradient descent (SGD) and prove sample complexity and runtime guarantees beyond NTK~\citep{damian2022neural, abbe2022merged, mousavi2022neural, li2020learning, abbe2021staircase,allen2019learning,allen2019can, xu2023over, wu2023learning, daniely2020learning, allen2020backward, suzukibenefit, telgarsky2022feature, chen2020towards, nichani2022escapeNTK}. These modified algorithms often contain multiple stages that optimize different blocks of parameters and/or use different learning rates or regularization strengths. For example, the work of~\citet{li2020learning} uses a non-standard parameterization for two-layer neural networks and runs a two-stage algorithm with sharply changing gradient clipping strength; \citet{damian2022neural} use one step of gradient descent with a large learning rate and then optimize only the last layer of the neural net. \citet{nichani2022escapeNTK} construct non-standard regularizers based on the NTK feature covariance matrix, in order to prevent the movement of weights in certain directions which hinder generalization. Additionally, \citet{abbe2022merged} only train the hidden layer for $O(1)$ time in the first stage of their algorithm, and then train the second layer to convergence in the second stage. They do not study how the population loss decreases during the first stage. However, oftentimes vanilla (stochastic) gradient descent with a constant learning rate empirically converges to a minimum with good generalization error. Thus, the modifications are arguably artifacts tailored to the analysis, and to some extent, over-using the modification may obscure the true power of gradient descent.

Another technique to analyze optimization dynamics uses the mean-field approach~\citep{chizat2018global, mei2018mean, mei2019mean,sirignano2018mean, dou2021training,abbe2022merged, javanmard2020analysis,sirignano2020mean, wei2019regularization}, which views the collection of weight vectors as a (discrete) distribution over $\R^d$ (where $d$ is the input dimension) and approximates its evolution by an infinite-width neural network, where the weight vector distribution is continuous and evolves according to a partial differential equation. 
However, these works do not provide an end-to-end polynomial runtime bound on the convergence to a global minimum in the concrete setting of two-layer neural networks (without modifying gradient descent). For example, the works of~\citet{chizat2018global} and~\citet{mei2018mean} do not provide a concrete bound on the number of iterations needed for convergence. \citet{mei2019mean} provide a coupling between the trajectories of finite and infinite-width networks with exponential growth of coupling error but do not apply it to a concrete setting to obtain a global convergence result with sample complexity guarantees. (See more discussion below and in Section~\ref{sec:main_results}.) \citet{wei2019regularization} achieve a polynomial number of iterations but require exponential width and an artificially added noise. In other words, these works, without modifying gradient descent, cannot prove convergence to a global minimizer with polynomial width and iterations. 

In this paper, we provide a mean-field analysis of projected gradient flow on two-layer neural networks with \textit{polynomial width} and quartic activations. Under a simple data distribution, we demonstrate that the network converges to a non-trivial error in \textit{polynomial time} with \textit{polynomial samples}. Notably, our results show a sample complexity that is superior to the NTK approach.


Concretely, the neural network is assumed to have unit-norm weight vectors and no bias, and the second-layer weights are all $\frac{1}{m}$ where $m$ is the width of the neural network. The data distribution is uniform over the sphere. The target function is of the form 	$y(x) = h(\truevector^\top x)$ where $h$ is an \textit{unknown} quartic link function and $\truevector$ is an unknown unit vector. 
Our main result (\Cref{thm:empirical-main}) states that with $n = O(d^{3.1})$ samples,  a polynomial-width neural network with random initialization converges in polynomial time to a non-trivial error, which is statistically not achievable by any kernel method with an inner product kernel using $n \ll d^4$ samples.  To the best of our knowledge, our result is the first to demonstrate the advantage of \textit{unmodified} gradient descent on neural networks over kernel methods.

The rank-one structure in the target function, also known as the single-index model, has been well-studied in the context of neural networks as a simplified case to demonstrate that neural networks can learn a latent feature direction $\truevector$ better than kernel methods~\citep{mousavi2022neural, abbe2022merged, damian2022neural, bietti2022learning}. 
Many works on single-index models study (stochastic) gradient descent in the setting where only a single vector (or ``neuron'') in $\mathbb{R}^d$ is trained. This includes earlier works where the link function is monotonic, or the convergence is analyzed with quasi-convexity ~\citep{kakade2011efficient,hazan2015beyond,mei2018landscape,soltanolkotabi2017learning}, along with more recent work~\citep{tan2019online, vardi2021learning, arous2021online} for more general link functions. Since these works only train a single neuron, they have limited expressivity in comparison to a neural network, and can only achieve zero loss when the link function equals the activation. We stress that in our setting, the link function $h$ is unknown, not necessarily monotonic, and does not need to be equal to the activation function. We show that in this setting, the first layer weights will converge to a \textit{distribution} of neurons that are correlated with but not exactly equal to $\truevector$, so that even without bias terms, their mixture can represent the link function $h$. Our analysis demonstrates that gradient flow, and gradient descent with a consistent inverse-polynomial learning rate, can \em simultaneously \em learn the feature $\truevector$ and the link function $h$, which is a key challenge that is side-stepped in previous works on neural-networks which use two-stage algorithms~\citep{mousavi2022neural, abbe2021staircase, abbe2022merged, damian2022neural, barak2022hidden}. The main novelty of our population dynamics analysis is designing a potential function that shows that the iterate stays away from the saddle points.

Our sample complexity results leverage a coupling between the dynamics on the empirical loss for a finite-width neural network and the dynamics on the population loss for an infinite-width neural network. The main challenge stems from the fact that some exponential coupling error growth is inevitable over a certain period of time when the dynamics resemble a power method update. 
Heavily inspired by~\citet{li2020learning}, we address this challenge by using a direct and sharp comparison between the growth of the coupling error and the growth of the signal.
In contrast, a simple exponential growth bound on the coupling error similar to the bound of~\citet{mei2019generalization} would result in a $d^{O(1)}$ sample complexity, which is not sufficient to outperform NTK.

As a corollary of our results for projected gradient flow, we show that projected gradient descent with polynomially small step size and polynomially many iterations can converge to a low error, again with sample complexity that is superior to NTK. Our proof is reminiscent of the standard bound on the error of Euler's method --- since our projected gradient flow runs for $O(\log d)$ time, a $\frac{1}{\poly(d)}$ learning rate suffices for the projected gradient descent trajectory to remain close to the projected gradient flow trajectory.